% Nouveau EGA V (cote 25): archive cote 25, site page 73 (PDF page 74).
% typescript. Draft compared: none.
% AI-transcribed by Claude Opus 5.5 (claude-opus-5-5, adaptive thinking, effort high), Message Batches API, batch ega25b (msgbatch_01A7n5YuXubUPNK1zsepi8RG), 2026-11-13; not yet checked by eye.
\scanpage{25/73}
\mspage{68}

relati\var{\xout{v}}{}{typewriter strike-out}{ega25_25_073-v0001}vement  au s chéma de base S:  V$^1$- V"  est en\var{\xout{xx}}{}{typewriter strike-out}{ega25_25_073-v0002} effet lisse sur S
(l'étant sur X)  de dimension relative égale à \uline{une} de moins que celle de
Grass$^m$ sur S  (comme on voit en faisant k=1 dans la formule plus haut).
Ainsi  les résultats d\unsure{e}  loc cit  s'appliquent, on trouve en particulier
le fait que l'ensemble des points  supersinguliers  de $\phi^m$ rel à Grass$^m$
n'est autre que V" \hand{$\cup$} V$^2$ , où  V$^2$ est le sous-\var{\inserted{pré}}{pré}{insertion}{ega25_25_073-v0003}schéma de ramification
de  V$^1$-V" \hand{$\longrightarrow$} G\unsure{r}ass$^m$ .  On peut donc dire encore que, en dehors de V", les
zéros supersinguliers résultent du collapsage de (au moins) deux zéros
singuliers ordinaires (mais on  n'est pas obligé de le dire\hand{)}).

\var{\margin{$|$}}{\placedmargin{$|$}}{marginal note placed}{ega25_25_073-v0004} De cette façon, on a essentiellement l'équivalent de 8.7. a) et b). Il \var{\margin{$+$}}{\placedmargin{$+$}}{marginal note placed}{ega25_25_073-v0005}
doit être possible de donner un équivalent pour  8.7. c), en \var{\xout{xxx}}{}{typewriter strike-out}{ega25_25_073-v0006} utilisant
la \var{\xout{structure}}{}{typewriter strike-out}{ega25_25_073-v0007} description explicite du fibré tangent  à Grass$^m$ (analogue à
celle du cas m=1); \var{\deleted{et qui}}{}{deletion}{ega25_25_073-v0008} \var{\inserted{cela}}{cela}{insertion}{ega25_25_073-v0009} implique \var{\inserted{encore}}{encore}{insertion}{ega25_25_073-v0010} que  pour  un point \var{\inserted{géométrique}}{géométrique}{insertion}{ega25_25_073-v0011} de  V$^1$-V"  non
ramifié au dessus  de Grass$^m$, la connaissance de son image dans Grass$^m$
implique celle de son image dans P, pourvu que la première image soit un
point lisse de l'image fermée \var{\inserted{de}}{de}{insertion}{ega25_25_073-v0012} V$^1$ dans Grass$^m$  (on suppose S = spectre d'
un corps).  Je peux fournir un énoncé \var{\xout{xxxxxx}}{}{typewriter strike-out}{ega25_25_073-v0013} plus précis sur demande.

Une fois ce point acquis,  on a des corollaires évidents généralisant
8.8. , 8.9. , 8.11.   Il est sans doute tout aussi possible d'énoncer dans
le cas m quelconque  les autres propositions du par.8. Si cela demande des
efforts de rédaction  supplémentaires, on peut renoncer à ce\unsure{tte} généralisa-%
tion\hand{,} \var{\xout{mêm}}{}{typewriter strike-out}{ega25_25_073-v0014} même si on inclut les développements différentiels précédents.
Il en est de même des résultats du n$^o$ 9.

Quant au n$^o$ 10 ,  la situation qui y est étudiée se généralise au
cas m quelconque de la façon suivante. On se fixe une sous-variété linéai-
re C de P  de codimension m+1, et on considère l'espace projectif Q des
sous-v\unsure{ariét}és linéaires L de P de codimension m passant par C. Q est une

\var{\margin{73}}{\placedmargin{73}}{marginal note placed}{ega25_25_073-v0015}
